\documentclass[11pt,a4paper]{article}

% =========================================================
% PACKAGES
% =========================================================
\usepackage[T1]{fontenc}
\usepackage[utf8]{inputenc}
\usepackage{lmodern}

\usepackage{amsmath,amssymb,amsfonts}
\usepackage{physics}
\usepackage{bm}
\usepackage{siunitx}

\usepackage{graphicx}
\usepackage{float}
\usepackage{subcaption}

\usepackage{booktabs}
\usepackage{array}

\usepackage{hyperref}
\hypersetup{
    colorlinks=true,
    linkcolor=blue,
    citecolor=blue,
    urlcolor=blue
}

\usepackage[
    margin=2.5cm
]{geometry}

% =========================================================
% TITLE INFORMATION
% =========================================================
\title{\textbf{Effects of Foreground Mitigation Placement in the HERA Analysis Pipeline on 21 cm Signal Recovery.}}

\author{
    Bernard Weich\\
    Department of Physics, Stellenbosch University\\
    Stellenbosch, South Africa\\
    \texttt{27068811@sun.ac.za}
}

\date{\today}

% =========================================================
% DOCUMENT
% =========================================================
\begin{document}

\maketitle

\begin{abstract}
A concise summary of the motivation, methods,
main results, and conclusions of the work.
Typically 150--250 words.
\end{abstract}

\section{Introduction}

Provide background, motivation, literature review,
and clearly state the objectives of the work.

HERA is a radio interferometer located at the SARAO site in the Karoo area in South Africa. The setup consists of 350 parabolic dishes, each with an antenna feed, 14 m in diameter, and centers 14.6 m apart. The core of the setup is made up of 320 of these dishes. Furthermore, two concentric hexagonal rings follow around the core made up of 30 dishes total.\\

HERA is excellent for studying large-scale structures due to large redundancy in baselines with small centre-to-centre distances. This large redundancy produces a better signal-to-noise ratio than for small structures. The maximum baseline length of the setup is 876 m, and the minimum is 14.6 m. Thus, HERA spans a wide range of cosmological scales that it can probe.\\

In the first results of the phase two implementation, after detailed analysis on a 14-night subset of data, constraints on the cosmological power spectrum with \(2\sigma\) confidence were determined. It was found that the constraints were stronger for redshift near EOR (Epoch of Reionization), and less constraining near the Cosmic Dawn. The paper also reports increased frequency bandwidth to 50-250 MHz, and that the main systematic is mutual coupling which leaks foreground power into low-k modes \((k=0.35 \longrightarrow k=0.55 \text{hMPc}^{-1}\).\\

So far, HERA Phase 2 has not been able to recover the faint 21cm signal due to this leakage. The focus of this paper is to investigate different configurations of the analysis pipeline used in the first results paper to try to reduce foreground leakage into these low k-modes.

\section{Theory}

Present the theoretical framework.

\begin{equation}
    H = \frac{p^2}{2m} + V(x)
\end{equation}

Define all symbols when first introduced.

\section{Experimental Methods}
% or Numerical Methods / Simulation Methods

Describe the experimental setup, materials,
equipment, computational methods, or simulations.

\section{Results}

Present results with figures and tables.

\begin{figure}[H]
    \centering
    \includegraphics[width=0.7\textwidth]{example_figure}
    \caption{Example figure caption.}
    \label{fig:example}
\end{figure}

\begin{table}[H]
    \centering
    \caption{Example table.}
    \begin{tabular}{cc}
        \toprule
        Parameter & Value \\
        \midrule
        A & 1.23 \\
        B & 4.56 \\
        \bottomrule
    \end{tabular}
\end{table}

\section{Discussion}

Interpret the results, discuss uncertainties,
limitations, and compare with previous work.

\section{Conclusion}

Summarize the main findings and suggest future work.

\section*{Acknowledgements}

Acknowledge funding, supervisors,
and collaborators.

% =========================================================
% REFERENCES
% =========================================================
\bibliographystyle{unsrt}
\bibliography{references}

\end{document}
```
